\section{Finite execution and the recorded certificate}
\label{sec:execution}

The computation used below was completed on 4 October 2026 (UTC). This
section specifies its finite inputs, the tests for complete execution,
and the implication from those tests to nonexistence. The geometric and
arithmetic justification of each rejection is given in
Sections~\ref{sec:seed}--\ref{sec:exclusion} and
Appendix~\ref{sec:rounding}. A successful exit by itself is not the
certificate: the inputs, record counts, saved failures, and coverage of
all shards must agree as described here.

\subsection{Programs, records, and machine arithmetic}
\label{execute:programs}

The accompanying \texttt{verification/code/} directory contains the seven unchanged
source files \texttt{seed.cpp}, \texttt{covout.cpp}, \texttt{ub.cpp},
\texttt{patchx.cpp}, \texttt{atlas.cpp}, \texttt{stageB.cpp}, and
\texttt{stageC.cpp}. Their byte lengths and full SHA-256 digests are
listed in \texttt{verification/code/extraction.json}. The source hashes were checked
before the recorded run, and the packaged copies match those same bytes.
Only four translation units are
built as entry points:
\begin{center}
\begin{tabular}{lll}
\toprule
Executable & Translation unit & Role\\
\midrule
\texttt{covrun} & \texttt{covout.cpp} & Nine-phase seed cover\\
\texttt{arun} & \texttt{atlas.cpp} & Completion search and coarse atlas\\
\texttt{brun} & \texttt{stageB.cpp} & Certified refinement\\
\texttt{crun} & \texttt{stageC.cpp} & Vector cover and graph exclusion\\
\bottomrule
\end{tabular}
\end{center}
The other files are included through the macros in these translation
units. Their standalone sample drivers are not executed. In particular,
A uses \texttt{Asearch}, and C uses graph search mode~1.

A serialized \texttt{Record} contains a certified \texttt{Root} and a
radius $r$. Its interpretation is the set of exact Hadamard solutions
represented, up to the equivalences of Section~\ref{sec:seed}, by
$e(p+h)$ with horizontal $h$ and $\|h\|\le r$. The stored $p$ need not
itself be an exact solution. The Root data justify the normal-displacement
inequalities on this set; their construction is part of
Section~\ref{sec:phase-cover}.

The binary format is explicit. A file begins with an eight-byte unsigned
record count. Each record consists, in order, of a four-byte signed
integer $k$, 36 binary64 entries of $p$, $k$ rows of 36 binary64 entries
of $T$, and the eight binary64 values
\[
 C,\ n_0,\ \mathrm{id},\ \mathrm{ort},\ \mathrm{hor},\
 \mathrm{mnorm},\ \mathrm{froot},\ r.
\]
Thus a record occupies $356+288k$ bytes, and a valid zero-count file is
exactly eight bytes. The platform is little endian. No phase data are
converted to decimal between stages. The evidence checks parse every
record, validate finite bounded fields, and require the declared count
to reach the exact end of file. They also check that every saved C
failure is an unchanged, ordered subsequence of its assigned inputs.

The recorded platform was x86-64, Intel Xeon Platinum 8573C, Ubuntu
GCC 13.3.0, and glibc 2.39. The builds, prepared on the same UTC date
before the run, were
\begin{lstlisting}
g++ -O3 -march=native -ffp-contract=off covout.cpp -o covrun
g++ -std=gnu++17 -O3 -march=native -ffp-contract=off atlas.cpp -o arun
g++ -std=gnu++17 -O3 -march=native -ffp-contract=off stageB.cpp -o brun
g++ -std=gnu++17 -O3 -march=native -ffp-contract=off stageC.cpp -o crun
\end{lstlisting}
The first command used this compiler's default GNU C++17 dialect.
The arithmetic assumptions are binary64, round to nearest, separate
operations without contraction or fast-math reassociation, correctly
rounded binary64 square root, masked floating-point exceptions, and
neither flush-to-zero nor denormals-are-zero. Separate probes compiled
with the cover and A/B/C flag sets both passed and recorded
\texttt{FLT\_EVAL\_METHOD=0}, 32-bit \texttt{int}, 64-bit
\texttt{size\_t}, 64-bit \texttt{uint64\_t} equal to
\texttt{unsigned long}, eight-bit bytes, and MXCSR \texttt{0x1f80}.
Both probes recorded nearest rounding on entry and during the check,
and passed the square-root and gradual-underflow samples. These probes
and the program checks provide limited evidence for the stated platform
conditions, not a proof of all primitive-operation or compiler semantics.
Those semantics remain assumptions; Appendix~\ref{sec:rounding} supplies
the numerical bounds under them. The retained
\texttt{build/build-manifest.json} identifies the exact sources,
binaries, compiler commands, probe outputs, and arithmetic evidence.

One additional code-generation condition is needed by
\texttt{coverSeeds::run}: both \texttt{pow(\ldots,2)} expressions must
be lowered to the binary64 multiplications analyzed in
Appendix~\ref{sec:rounding}. For the actual \texttt{covrun}, inspection
of the saved disassembly finds the scale multiplication, its square,
subtraction, square root, addition, and second square at hexadecimal
addresses \texttt{6035}, \texttt{6039}, \texttt{603d},
\texttt{605c}, \texttt{606d}, and \texttt{608b}, respectively. There
is no dynamic \texttt{pow} symbol. Absence of that symbol alone would
not establish the required instruction sequence. The recorded binary is
identified by its full SHA-256 digest, so this check applies to that build.
The specified lowerings address only this code-generation condition, not
the full arithmetic model itself. A new build requires
this check again. Its binary hash, like any native build hash, need not
equal that of the recorded binary.

\subsection{Specified finite control flow}
\label{execute:control}

All invocations use the C numeric locale and fresh output paths.
The cover is run to normal completion before A reads its output.
A uses all four parts $i=0,1,2,3$ with \texttt{parts=4}; part $i$
processes exactly the center indices congruent to $i$ modulo~4.
Its constants are seed radius \texttt{.0320001}, local target
\texttt{.103}, global target \texttt{.125},
\texttt{levels=3}, \texttt{iters=2}, and \texttt{first=1}.

The following are the program invocations, with $i$ and $j$ standing for
the indicated shard and batch indices. The accompanying
\texttt{verification/code/proof\_run.py} runner was used in this execution and
performs the validation between stages; this display only specifies the
program arguments. A separate bounded supervisor retained the runner's
exit and process-cleanup record.
\begin{lstlisting}
export LC_ALL=C
./covrun centers.txt
./arun centers.txt i 4 coarsei.bin badi.txt
./brun .07 .013 mid.bin midbad.bin - \
  coarse0.bin coarse1.bin coarse2.bin coarse3.bin
./brun .016 .005 fine.bin finebad.bin batches/f mid.bin
./crun 12 4 1 .042 1 1 5000000 0 1 \
  batches/badj.bin batches/fj.bin
./brun .011 .0038 fb1.bin fb1bad.bin - ALL_NONEMPTY_C_BAD_FILES
./crun 12 4 1 .042 1 1 5000000 i 4 fb1cbadi.bin fb1.bin
\end{lstlisting}
There is one initial C invocation for every batch named by the completed
\texttt{batches/fdone.txt} marker. In C the first seven parameters mean
initial grid $12$, four layers, one linear-support iteration, proposed
cluster radius $.042$, detail~1, graph mode~1, and graph-step cap
$5{,}000{,}000$. Thus the single-vector grids are $12,24,48,96$.
The unchanged default number of clustering passes is two.

The mid and fine B invocations are sequential. A batching B invocation
writes consecutive groups of at most 256 atlas records, replacing each
batch radius by its target $.016$, and atomically renames each completed
batch. The checked concatenation of all batch records equals
\texttt{fine.bin}, except for that radius replacement. Every original
fine radius is at most $.016$, so the replacement enlarges the represented
sets. The completion marker agrees with both the number of batches and
the sum of their record counts. The recorded run waited for B fine to
finish before starting C; it used up to sixteen independent C workers.
Concurrency changes scheduling only: workers have distinct output files
and do not share mutable search state.

For every initial C process, the normal status is 0 exactly when its
saved bad count is zero, and 1 exactly when that count is positive.
The four reported failure categories---uncapped graph candidate,
graph-step cap, bin overflow, and clustering failure---sum to that bad
count. None of them is an exclusion. Every nonempty-count bad file is
passed to B fallback in C-locale filename order, including failures whose
locations differ from those in another run. Final C uses all four
modulo-4 parts of the fallback output. If initial C has no bad records,
fallback and final C are unnecessary; that shortcut was not used in the
recorded run.

The required terminal conditions are:
\begin{enumerate}
\item successful builds, a normally completed cover, and a complete
      center file;
\item four A exits 0, all assigned centers processed, terminal
      \texttt{failedBalls=0}, and four empty bad text files;
\item both B exits 0, all input parents processed, agreement of logged
      and saved output counts, and zero-count B bad files;
\item every initial C batch processed exactly once with consistent
      exit, done count, bad count, and saved records;
\item every initial C bad record included in fallback, fallback exit 0
      with all parents processed and zero saved failures, and four
      final C exits 0 with total done equal to the fallback count and
      zero saved failures.
\end{enumerate}
Missing files, truncated data, an unexpected exit, resource exhaustion,
an unresolved A or B parent, or a final C bad record prevent acceptance.
Recoverable search diagnostics are distinguished from these terminal
conditions. For example, A's \texttt{ownfail} counts unsuccessful
proposals of a local Root cover that are followed by subdivision. A0
reported such a diagnostic in this run; its terminal unresolved count
was nevertheless zero.

\subsection{Finite bounds and conservative stopping}
\label{execute:bounds}

The control flow has explicit finite bounds. Before filtering, the
nondecreasing five-bin lists for the seed cover number at most
\[
 \binom{116}{5}=160{,}389{,}488,
\]
and for any fixed shared bin the four ordered tail bins number at most
\[
 \binom{115}{4}=6{,}913{,}340.
\]
The former bound is below the signed 32-bit limit. Pair traversals use
range loops and 64-bit counters. All bin differences lie between $-111$
and $111$; their nine-coordinate square sum is at most $9\cdot111^2$,
so the integer distance and bucket arithmetic fit the stated types.
The stored five-coordinate differences also fit the signed eight-bit
type used by the cover.

A retains nonfinal unitary grid lists only through grid~32; the grid~64
has at most $64^4/8=2{,}097{,}152$ cell tests per seed. Further tests
are bounded by the $1{,}000{,}000$ extra-node cap, and the deepest
grid is 512. Reaching the cap or reaching the deepest grid without a
certificate records an unresolved seed. Trial-direction and Root
proposal failures can lose an opportunity to cover a cell, but cannot
make an uncovered cell successful. The local candidate list uses
\texttt{topK=4} in an array of size~8.

The Root dimension is at most $k=10$. Atlas sizes and individual record
file counts are capped at $2{,}000{,}000$, and concatenated B/C input
counts at $8{,}000{,}000$. For B the validated widths and radii give
$L\le1262$ and hence axis indices at most 2524; the product grid is
checked against $50{,}000{,}000$ before multiplication. Mixed-radix
strides and marks therefore fit the stated index types. Each unresolved
box is split to depth at most six, and the $2^k$ child masks fit an
integer. Only successfully completed parents are inserted into B's
list of earlier covered parents. A grid-cap rejection is a parent
failure, not a completed parent.

C retains at most $2{,}000{,}000$ phase bins. The first grid has
$12^5=248{,}832$ cells; a later grid tests at most 32 children per
retained bin. Clustering permits at most 2048 balls. Its graph bitsets
have 32 words of 64 bits, and all members have indices below the current
node count. Word shifts are unsigned. The graph pair counts fit an
integer, and the joint test has at most 14 rows. Recursive clique
search decreases the requested clique size or removes a candidate;
the step cap supplies an additional bound. A cap sets the conservative
candidate/give-up flag and saves the record. It does not return a
negative certificate. Consequently no finite resource cap in the
specified run can establish nonexistence by silently dropping a region.

\subsection{The finite-execution implication}
\label{execute:implication}

\begin{theorem}[Finite-execution certificate]
\label{execute:finite-theorem}
Use the seven specified source files and the arithmetic conditions of
Appendix~\ref{sec:rounding}. Suppose the invocations in
Section~\ref{execute:control} meet all of its terminal conditions, with
the record provenance and batch checks stated there. Then four mutually
unbiased orthonormal bases in $\C^6$ do not exist.
\end{theorem}

\begin{proof}
Suppose four such bases exist. Normalize the first to the standard basis
and write the second as $H/\sqrt6$. Proposition~\ref{seed:cover} places
the seed coordinates of an equivalent exact Hadamard matrix in one of
the generated seed balls. All A shards together process that seed.
Proposition~\ref{phase:A-cover} therefore places an equivalent exact
solution in a coarse Record. Each equivalence preserves extendibility
by two further bases.

Apply Proposition~\ref{refine:cover-preservation} first to B mid and
then to B fine, and Corollary~\ref{refine:batch-cover} to the emitted
batches. They place
an equivalent exact solution in a record of \texttt{fine.bin}, and
therefore in its checked, possibly enlarged batch record. Completeness
of the marker and the per-batch done counts ensures that initial C
processes that record. A negative C result would contradict its
extendibility by Proposition~\ref{exclude:local-certificate}. Hence it
must occur among the saved initial failures. If there are none, this
already gives a contradiction. Otherwise the fallback branch is executed,
and the subsequence and fallback-input checks include that very record
in B fallback.

The successful fallback again covers its exact solutions, up to the
same equivalences, by the records of \texttt{fb1.bin}. The final four
modulo-4 shards partition those records without omission. Every final
record has a negative C result, contradicting extendibility. Each
coverage implication used here is justified by its geometric lemma and
the one-sided arithmetic bounds; counts and hashes establish the
identity and completeness of the finite chain, not those inequalities
in isolation.
\end{proof}

\subsection{The completed execution}
\label{execute:observed}

The execution began at 03:10:26 and completed at 04:24:56 on
4 October 2026 (UTC). Both the runner and its supervisor exited 0;
the supervisor recorded no cleanup or finalization errors, no remaining
children, and unchanged input hashes. The execution evidence includes
the raw command and stream records, exact file contents, record provenance,
batch concatenation, actual fallback inputs, and final modulo partition.
Together with the validation records, these account for all 100 native
invocations and show no final unresolved records. They document the finite
execution, not an automatic verification of the geometric or arithmetic
lemmas.

Table~\ref{execute:counts} records this computation. All four A shards exited 0,
with empty bad text files. Each B invocation exited 0 with an eight-byte
zero-count bad file. Initial C had 73 exits 0 and 15 exits 1, all
accounted for by saved failures. All final C shards exited 0 with
eight-byte zero-count bad files.

\begin{table}[htbp]
\centering
\small
\begin{tabular}{lrrl}
\toprule
Stage & Inputs processed & Output records & Unresolved records\\
\midrule
Seed cover & $151{,}744{,}086$ pairs & $249{,}878$ centers & ---\\
A0 & $62{,}470$ & 764 & 0\\
A1 & $62{,}470$ & 774 & 0\\
A2 & $62{,}469$ & 757 & 0\\
A3 & $62{,}469$ & 788 & 0\\
B mid & 3083 & 600 & 0\\
B fine & 600 & $22{,}379$ & 0\\
Initial C & $22{,}379$ & --- & 163 (all graph caps)\\
B fallback & 163 & 1127 & 0\\
Final C0 & 282 & --- & 0\\
Final C1 & 282 & --- & 0\\
Final C2 & 282 & --- & 0\\
Final C3 & 281 & --- & 0\\
\bottomrule
\end{tabular}
\caption{Observed execution on 4 October 2026 (UTC). The seed cover retained
22,829 five-bin $x$ lists. The fine records form 88 batches:
$87\cdot256+107=22{,}379$.}
\label{execute:counts}
\end{table}

The nonzero initial C failure files, in the actual order supplied to
B fallback, were
\begin{lstlisting}
bad10.bin bad3.bin bad4.bin bad5.bin bad52.bin bad53.bin
bad55.bin bad56.bin bad6.bin bad7.bin bad78.bin bad79.bin
bad8.bin bad80.bin bad84.bin
\end{lstlisting}
They contain all 163 unresolved records. There were no initial bin
overflows, clustering failures, or uncapped graph candidates. There
were no final failures of any category. Thus all 100 program invocations
in the full pipeline are accounted for: one cover, four A processes,
two B refinements, 88 initial C processes, one B fallback, and four
final C processes. The fifteen initial C exits 1 are resolved by that
fallback chain. No rejection threshold, graph cap, refinement width,
or mathematical source file was altered, and no additional refinement
level was inserted.

For scale, the wrapper-measured stage wall times, including polling and
validation, were 621.59 seconds for the cover, 2020.32 seconds for the
four parallel A shards, 1071.49 and 407.45 seconds for B mid and B fine,
and 300.23, 8.12, and 36.27 seconds for initial C, fallback B, and final C.
The outer supervisor measured 4470.33 seconds for the complete pipeline.
These are observations on this platform, not native CPU times or
running-time guarantees. No stage hit its wall deadline, and no retry
or resource-limit increase was used.

\begin{longtable}{@{}p{.19\textwidth}p{.77\textwidth}@{}}
\caption{SHA-256 digests of the actual principal data files. Full
per-file inventories accompany the computation archive.}
\label{execute:hashes}\\
\toprule
File & SHA-256\\
\midrule
\endfirsthead
\toprule
File & SHA-256\\
\midrule
\endhead
\texttt{centers.txt} & {\scriptsize\texttt{c48bc9a2be1eb2cdeea82ea1b030d7cb9f00a6c9563b63f0248dd115748a5789}}\\
\texttt{coarse0.bin} & {\scriptsize\texttt{77b0cda59484ffde0654db8d6008bcaab1daf1deefd00d4dc085d3d751418cd4}}\\
\texttt{coarse1.bin} & {\scriptsize\texttt{48ca7775dba57a87a10d154da8b172074b29c9a65b32e44ee169f23a48d24b8a}}\\
\texttt{coarse2.bin} & {\scriptsize\texttt{418b6d9cf338187c63635d26050940a13177515050a057b2bc67fd11a163e43e}}\\
\texttt{coarse3.bin} & {\scriptsize\texttt{8c23427e63e31cd18abbb80838e3289e646fd108fa850687ca7b857471376535}}\\
\texttt{mid.bin} & {\scriptsize\texttt{e97c74dde90a0e758df2cb5d9708f5bd6bd660ed920b3299f5df92ff4de11be4}}\\
\texttt{fine.bin} & {\scriptsize\texttt{9b5b743b61907277fa6687a0ad45677bbcd0e4c9bd44a0d93d258253fb52923d}}\\
\texttt{fb1.bin} & {\scriptsize\texttt{ef85ec73217ecb0c770d3c775396df88b60ca06918f77310df5dff2f6d9bf215}}\\
\bottomrule
\end{longtable}

The execution archive identified by
\texttt{verification/\allowbreak computation/README.md} preserves the numerical data and records
of this complete run, with internal path and host fields redacted: all sources
and binaries, center and record files, all 88 input batches and their C
bad files, fallback data, all final bad files, stdout and stderr,
commands, exits, platform information, and validation records. In the
execution archive, \texttt{full-run-01/pipeline/status.json} is the
whole-run record; \texttt{*.stage.json} files describe stages and
individual \texttt{*.status.json} files describe child processes.
The outer record is \texttt{full-run-01/run.json}.
\texttt{verification/REPRODUCE.md} gives rebuild and
execution instructions. Archived native binaries document this run;
they are not a portability claim.

The proof here uses only this newly completed chain and the values in
Tables~\ref{execute:counts} and~\ref{execute:hashes}. The complete
evidence packet for the execution reported on 7 September 2026 was not
recovered. Although the new principal data hashes and stage counts agree
with those reported for that run, this agreement does not authenticate
the earlier execution. Neither its reported results nor earlier timed
or sampled diagnostics substitute for any stage of the new full run.
